Two plate capacitors 1 and 2 are charged and then disconnected from the source, so the charges on their plates stay. In both, the upper plate is positive. Between the plates of capacitor 1 is an electron, between those of capacitor 2 an alpha particle.

\begin{center}
\begin{tabular}{lcc}
 & capacitor 1 & capacitor 2 \\ \hline
charge on the plates & $Q$ & $2Q$ \\
area of each plate & $A$ & $A/2$ \\
distance between the plates & $d$ & $2d$ \\
particle between the plates & electron ($-e$) & alpha particle ($+2e$)
\end{tabular}
\end{center}

\begin{abcliste}
\abc By what factor is the field in capacitor 2 stronger than in capacitor 1?
\abc By what factor is the electric force on the particle in capacitor 2 larger than on the particle in capacitor 1?
\abc Do the two forces point the same way or opposite ways?
\abc Now both capacitors stay connected to voltage sources instead, capacitor 2 at twice the voltage of capacitor 1 (areas and distances as in the table). Answer (a) and (b) again.
\end{abcliste}
